
\documentclass[a4paper,10pt]{report}

%\usepackage[landscape]{geometry}
\usepackage[cm-default]{fontspec}
\setromanfont{FreeSerif}
\setsansfont{FreeSans}
\setmonofont{FreeMono}
\usepackage[utf8x]{inputenc}
\usepackage{fontspec}
\usepackage{xunicode}
\usepackage{xltxtra}
\usepackage{xgreek}
\usepackage{amsmath}
\usepackage{unicode-math}
\usepackage{ulem}
\usepackage{color}
\usepackage{verbatim}
\usepackage{nopageno}
\usepackage{graphicx}
\usepackage{textpos}
\setlength{\TPHorizModule}{1cm}
\setlength{\TPVertModule}{1cm}
%\usepackage[colorgrid,texcoord]{eso-pic}
\usepackage[outline]{contour}
\usepackage{wrapfig}
\usepackage{url}
\usepackage{color}
\usepackage{tikz}
\usepackage{fancybox,fancyhdr}
\usepackage{subfigure}
\usepackage{pstricks}
\usepackage{epsfig}
\usepackage{multicol}
\usepackage{listings}
\usepackage{enumerate}
\usepackage{hyperref}
\hypersetup{
  bookmarks=true,
  bookmarksopen=true,
  pdfborder=false,
  pdfpagemode=UseNone,
  raiselinks=true,
  pdfhighlight={/P},
  colorlinks,
  citecolor=black,
  filecolor=black,
  linkcolor=black,
  urlcolor=black
}


\usepackage{fontspec}
\usepackage{xunicode}
\usepackage{xltxtra}


% Margins
%\setlength{\textwidth}{24cm}

%\setlength{\voffset}{0in}
%\setlength{\textheight}{6.8in}
% Colors
\definecolor{cbrown}{rgb}{0.49,0.24,0.07}
\definecolor{cmpez}{rgb}{0.92,0.88,0.79}
\definecolor{cmpez2}{rgb}{0.62,0.33,0.34}
\definecolor{cwhite}{rgb}{1,1,1}
\definecolor{cblack}{rgb}{0,0,0}
\definecolor{cred}{rgb}{0.9,0.15,0.15}
\definecolor{clblue}{rgb}{0.57,0.85,0.97}
\definecolor{corange}{rgb}{0.86,0.49,0.18}
\definecolor{clorange}{rgb}{0.95,0.84,0.61}
\definecolor{cyellow}{rgb}{1,0.95,0.16}
\definecolor{cgreen}{rgb}{0.69,0.8,0.31}
\definecolor{clbrown}{rgb}{0.78,0.67,0.43}
\definecolor{cmagenta}{rgb}{0.79,0.05,0.54}
\definecolor{cgray}{rgb}{0.76,0.73,0.63}
\definecolor{bluesite}{rgb}{0.05,0.43,0.67}
\definecolor{redsite}{rgb}{0.90,0.15,0.15}
\definecolor{backsite}{rgb}{0.07,0.07,0.07}
% ----------- from Costas -----------------------------

%---------------   Main Fonts    -----------------%   

% \setmainfont[Mapping=tex-text]{Calibri}
% \setmainfont[Mapping=tex-text]{Times New Roman}
%\setmainfont[Mapping=tex-text]{Droid Serif}
%\setmainfont[Mapping=tex-text]{Cambria}
%\setmainfont[Mapping=tex-text]{cm-unicode}
% \setmainfont[Mapping=tex-text]{Gentium}
% \setmainfont[Mapping=tex-text]{GFS Didot}
% \setmainfont[Mapping=tex-text]{Comic Sans MS}
 \setmainfont[Mapping=tex-text]{Ubuntu}
% \setmainfont[Mapping=tex-text]{Myriad Pro}
% \setmainfont[Mapping=tex-text]{CMU Concrete}
%\setmainfont[Mapping=tex-text]{DejaVu Sans}
%\setmainfont[Mapping=tex-text]{KerkisSans}
%\setmainfont[Mapping=tex-text]{KerkisCaligraphic}
% \setmainfont[Mapping=tex-text]{Segoe Print}
% \setmainfont[Mapping=tex-text]{Gabriola}

%------------------------------------------------%

% ------------ Mathematics Fonts ----------------%
\setmathfont{Asana-Math.ttf}
% -----------------------------------------------%
% Misc
\author{Κασωτάκη Ε. - Σμαραγδάκης Κ.}
\title{www.math24.gr}
% ---------------- Formation --------------------%


\setlength\topmargin{-0.7cm}
\addtolength\topmargin{-\headheight}
\addtolength\topmargin{-\headsep}
\setlength\textheight{26cm}
\setlength\oddsidemargin{-0.54cm}
\setlength\evensidemargin{-0.54cm}
%\setlength\marginparwidth{1.5in}
\setlength\textwidth{17cm}

\RequirePackage[avantgarde]{quotchap}
\renewcommand\chapterheadstartvskip{\vspace*{0\baselineskip}}
\RequirePackage[calcwidth]{titlesec}
\titleformat{\section}[hang]{\bfseries}
{\Large\thesection}{12pt}{\Large}[{\titlerule[0.9pt]}]
%--------------------------------------------------%
\usepackage{draftwatermark}
\SetWatermarkText{www.math24.gr}
\SetWatermarkLightness{0.9}
\SetWatermarkScale{0.6}
%--------------------------------------------------%


\usepackage{xcolor}
\usepackage{amsthm}
\usepackage{framed}
\usepackage{parskip}

\colorlet{shadecolor}{bluesite!20}

\newtheorem*{orismos}{Ορισμός}

\newenvironment{notation}
  {\begin{shaded}\begin{theorem}}
  {\end{theorem}\end{shaded}}




% Document begins
\begin{document}
\pagestyle{fancy}
\fancyhead{}
\fancyfoot{}
\renewcommand{\headrulewidth}{0pt}
\renewcommand{\footrulewidth}{0pt}


\fancyhead[LO,LE]{
 %\textblockcolor{backsite}
 %\begin{textblock}{5}(-2,-0.55)
  %\rule{0cm}{1cm}
 %\end{textblock}
 \textblockcolor{bluesite}
 \begin{textblock}{5}(-1.5,-0.55)
  \rule{0cm}{1cm}
 \end{textblock}
 %\textblockcolor{bluesite}
 %\begin{textblock}{14}(5.5,-0.55)
 % \rule{0cm}{1cm}
 %\end{textblock}
 \begin{textblock}{0}(-1,-0.25)
 \color{cwhite} \begin{Large}www.math24.gr\end{Large}
 \end{textblock}
\begin{textblock}{0}(16,-0.4)
 \color{cwhite} \includegraphics[height=1.5cm]{math24_logo.png}
 \end{textblock}
\textblockcolor{white}
\begin{textblock}{17}(-1,28)
 \color{backsite} \begin{small}Copyright \textcopyright 2011, Κασωτάκη Ε.(ikasotaki@gmail.com) - Σμαραγδάκης Κ.(kesmarag@gmail.com)\end{small}
 \end{textblock}
}
 
\begin{shaded}
\begin{center}
\huge \textbf{Φυλλάδιο Ασκήσεων}\\
%Πρόσθεση ρητών αριθμών
\end{center} 
\textbf{Μαθηματικά Α' Γυμνασίου} \hfill \textbf{Ημερομηνία Παράδοσης : \hspace{2em} }
\subsection*{Ονοματεπώνυμο :\hfill  \hspace{5em}}
\end{shaded}
\vspace{2em}
\begin{itemize}
 \item Διαίρεση 2 ομόσημων ρητών αριθμών
 \item Διαίρεση 2 ετερόσημων ρητών αριθμών
 \item Διαίρεση κλασμάτων
\end{itemize}
\section*{Θεωρία - Διαίρεση 2 ομόσημων ρητών αριθμών\hfill \small{}}
\begin{itemize}
 \item \textbf{Ομόσημοι} λέγονται οι αριθμοί που έχουν το ίδιο πρόσημο. 
 \item \underline{ Κανόνας για την διαίρεση 2 ομόσημων ρητών αριθμών:} \\
       Για να \textbf{διαιρέσουμε 2 ομόσημους} ρητούς αριθμούς, 
       \textbf{διαιρούμε} τις απόλυτες τιμές τους και στο πηλίκο βάζουμε το πρόσημο "\textbf{+}"\\
       Δηλαδή $+\cdot +=+$ και $-\cdot -=+$\\
       \textbf{π.χ} $24:3=+8$\\
       \textbf{π.χ} $27:2=13.5$\\
       \textbf{π.χ} $11.25:5=2.25$ \\
       \textbf{π.χ} $(-32):(-4)=+8$ \\
       \textbf{π.χ} $(-11.22):(-5.1)=+2.2$ \\
       \textbf{π.χ} $(-31.72):(-10)=3.172$ \\
\end{itemize}


\section*{Άσκηση 1  (Διαίρεση 2 ομόσημων ρητών αριθμών)\hfill \small{40 μονάδες}}
Να κάνετε τις παρακάτω διαιρέσεις:
\begin{multicols}{2}
\begin{enumerate}[1)]
 \item $1.5:3$
 \item $72.1:100$
 \item $81:3$
 \item $3:50$
 \item $2022:2$
 \item $(-0.45):(-3)$
 \item $(-21):(-7)$
 \item $(-380):(-5)$
 \item $(-565):(-5)$
 \item $(-0.1):(-2)$
 \item $333:9$
 \item $(-123):(-100)$
 \item $+23.4:(+2)$
 \item $(-725):(-5)$
 \item $702:3$
 \item $(-168):(-42)$
 \item $555:3$
 \item $(-81.5):(-10)$
 \item $1024:8$
 \item $(-111):(-3)$
\end{enumerate}
\end{multicols}

\newpage
\section*{Θεωρία - Διαίρεση 2 ετερόσημων ρητών αριθμών\hfill \small{}}
\begin{itemize}
 \item \textbf{Ετερόσημοι} λέγονται οι αριθμοί που έχουν διαφορετικό πρόσημο. 
 \item \underline{ Κανόνας για την διαίρεση 2 ετερόσημων ρητών αριθμών:} \\
       Για να \textbf{διαιρέσουμε 2 ετερόσημους} ρητούς αριθμούς, 
       \textbf{διαιρούμε} τις απόλυτες τιμές τους και στο πηλίκο βάζουμε το πρόσημο "\textbf{-}"\\
       Δηλαδή $-:+=-$ και $+ -=-$\\
       \textbf{π.χ} $(-16):(+4)=-4$\\
       \textbf{π.χ} $(-10.1):(+10)=-1.01$\\
       \textbf{π.χ} $(-128):(+8)=-16$ \\
       \textbf{π.χ} $(+11.7):(-3)=-3.9$ \\
       \textbf{π.χ} $(+162):(-9)=-18$ \\
       \textbf{π.χ} $(+1337):(-1000)=-1.337$ \\
\end{itemize}

\section*{Άσκηση 2  (Διαίρεση 2 ετερόσημων ρητών αριθμών)\hfill \small{40 μονάδες}}
Να κάνετε τις παρακάτω διαιρέσεις:
\begin{multicols}{2}
\begin{enumerate}[1)]
 \item $(-36):6$
 \item $(-36):12$
 \item $(-12.3):3$
 \item $(-53.5):5$
 \item $(-400):50$
 \item $13:(-0.1)$
 \item $24:(-0.2)$
 \item $39:(-10)$
 \item $6.6:(-2.2)$
 \item $720:(-90)$
 \item $(-3.9):3$
 \item $(52.5):(-7)$
 \item $(-2.7):(0.9)$
 \item $(0.12):(-3)$
 \item $(-23.1):(+100)$
 \item $(-0.75):(0.25)$
 \item $(-49.5):(5)$
 \item $164:(-82)$
 \item $(-113):(10)$
 \item $0.6:(-30)$
\end{enumerate}
\end{multicols}






\section*{Θεωρία - Διαίρεση ρητών αριθμών \hfill \small{}}
 $$ \frac{α}{β}=α\cdot \frac{1}{β} $$
 Για να διαιρέσουμε δύο ρητούς αριθμούς αρκεί να πολλαπλασιάσουμε το διαιρετέο με τον αντίστροφο του διαιρέτη 
\begin{itemize}
  \item      \textbf{π.χ} $\dfrac{2}{3}:\dfrac{5}{4}=\dfrac{2}{3}\cdot \dfrac{4}{5}=\dfrac{8}{15}$
  \item      \textbf{π.χ} $(-\dfrac{1}{3}):(-\dfrac{4}{5})=+\dfrac{1}{3}\cdot \dfrac{5}{4}=\dfrac{5}{12}$ 
  \item      \textbf{π.χ} $(-\dfrac{7}{10}):(+\dfrac{3}{5})=-(\dfrac{7}{10}\cdot \dfrac{5}{3})=-\dfrac{35}{30}=-\dfrac{7}{6}$ 
  \item      \textbf{π.χ} $\dfrac{1}{2}:(-\dfrac{9}{4})=-(\dfrac{1}{2}\cdot \dfrac{4}{9})=-\dfrac{2}{9}$  
\end{itemize}
\newpage
\section*{Άσκηση 3  (Διαίρεση Κλασμάτων)\hfill \small{20 μονάδες}}
Να κάνετε τις παρακάτω διαιρέσεις:
\begin{multicols}{2}
\begin{enumerate}[1)]
 \item $\dfrac{10}{3}:\dfrac{5}{4}$
 \item $(-\dfrac{8}{3}):(\dfrac{2}{7})$
 \item $(-\dfrac{9}{4}):(-\dfrac{3}{4})$
 \item $\dfrac{8}{5}:(-\dfrac{50}{40})$
 \item $(-\dfrac{11}{2}):(\dfrac{11}{2})$
 \item $(-\dfrac{1}{100}):(-\dfrac{60}{30})$
 \item $(+\dfrac{23}{3}):(-\dfrac{13}{9})$
 \item $\dfrac{1}{2}:\dfrac{2}{8}$
 \item $(-\dfrac{55}{3}):\dfrac{1}{125}$
 \item $(-\dfrac{1}{4}):(-\dfrac{1}{2})$
\end{enumerate}
\end{multicols}








\end{document}
